\documentclass[11pt,border=4pt]{standalone}
% Estilo comum dos esquemáticos do relatório (circuitikz).
% Uso: \documentclass[11pt,border=4pt]{standalone} e, em seguida, % Estilo comum dos esquemáticos do relatório (circuitikz).
% Uso: \documentclass[11pt,border=4pt]{standalone} e, em seguida, % Estilo comum dos esquemáticos do relatório (circuitikz).
% Uso: \documentclass[11pt,border=4pt]{standalone} e, em seguida, \input{estilo_circuitos}.
\usepackage[T1]{fontenc}
\usepackage[utf8]{inputenc}
\usepackage[brazilian]{babel}
\usepackage{lmodern}
\usepackage{amsmath}
\usepackage[american]{circuitikz}
\usetikzlibrary{calc,arrows.meta}

% Cores das figuras originais
\definecolor{anot}{RGB}{0,135,200}      % anotações: valores calculados, setas, resistências vistas
\definecolor{sombra}{RGB}{220,242,253}  % região do transistor nos modelos T e pi
\definecolor{term}{RGB}{62,124,148}     % nomes de terminais destacados (B, C, E)

% Fios com 0,4 pt e corpo dos componentes com 0,8 pt, como nas figuras originais
\ctikzset{bipoles/thickness=2}
.
\usepackage[T1]{fontenc}
\usepackage[utf8]{inputenc}
\usepackage[brazilian]{babel}
\usepackage{lmodern}
\usepackage{amsmath}
\usepackage[american]{circuitikz}
\usetikzlibrary{calc,arrows.meta}

% Cores das figuras originais
\definecolor{anot}{RGB}{0,135,200}      % anotações: valores calculados, setas, resistências vistas
\definecolor{sombra}{RGB}{220,242,253}  % região do transistor nos modelos T e pi
\definecolor{term}{RGB}{62,124,148}     % nomes de terminais destacados (B, C, E)

% Fios com 0,4 pt e corpo dos componentes com 0,8 pt, como nas figuras originais
\ctikzset{bipoles/thickness=2}
.
\usepackage[T1]{fontenc}
\usepackage[utf8]{inputenc}
\usepackage[brazilian]{babel}
\usepackage{lmodern}
\usepackage{amsmath}
\usepackage[american]{circuitikz}
\usetikzlibrary{calc,arrows.meta}

% Cores das figuras originais
\definecolor{anot}{RGB}{0,135,200}      % anotações: valores calculados, setas, resistências vistas
\definecolor{sombra}{RGB}{220,242,253}  % região do transistor nos modelos T e pi
\definecolor{term}{RGB}{62,124,148}     % nomes de terminais destacados (B, C, E)

% Fios com 0,4 pt e corpo dos componentes com 0,8 pt, como nas figuras originais
\ctikzset{bipoles/thickness=2}

% Amplificador diferencial com MOSFETs da Questão 3 (fig:q3-circuito).
% Compilar com pdflatex; o PDF resultante vai para figuras/fig_q3_circuito.pdf
% e entra no relatório em escala 1.

\begin{document}
\begin{circuitikz}
  % Geometria em cm: nó S na origem, colunas de dreno e fonte em x = -2,3 e x = 2,3.
  \def\xL{-2.3}
  \def\xR{2.3}
  \def\yF{0.5}   % terminais de fonte de M1 e M2
  \def\yD{3.3}   % nós D1 e D2
  \def\yT{5.1}   % trilho de V+
  \def\yB{-2.0}  % extremo inferior de RS (V-)
  \def\xO{4.1}   % terminal de saída vO2

  % Transistores NMOS com seta na fonte; M2 espelhado, com a porta à direita.
  \node[nmos, arrowmos, anchor=source] (M1) at (\xL,\yF) {};
  \node[nmos, arrowmos, xscale=-1, anchor=source] (M2) at (\xR,\yF) {};

  % Ramo superior: RD, nós D1 e D2, trilho de V+.
  \draw (M1.drain) -- (\xL,\yD) to[R] (\xL,\yT) -- (\xR,\yT) to[R] (\xR,\yD) -- (M2.drain);
  \draw (0,\yT) -- (0,\yT+0.4) node[vcc]{};
  % Saída vO2 no dreno de M2.
  \draw (\xR,\yD) to[short, -o] (\xO,\yD);
  % Entradas nas portas.
  \draw (M1.gate) to[short, -o] ++(-1.0,0) coordinate (V1);
  \draw (M2.gate) to[short, -o] ++(1.0,0) coordinate (V2);
  % Ramo inferior: fontes ligadas em S, RS até V-.
  \draw (M1.source) -- (\xL,0) -- (\xR,0) -- (M2.source);
  \draw (0,0) to[R] (0,\yB) node[vee]{};
  % Pontos de junção.
  \node[circ] at (0,\yT) {};
  \node[circ] at (\xL,\yD) {};
  \node[circ] at (\xR,\yD) {};
  \node[circ] at (0,0) {};

  % Sentidos de corrente (as coordenadas guardam a ponta de cada seta).
  \draw[-latex, thick] (\xL+0.3,\yD-0.32) -- ++(0,-0.65) coordinate (ID1);
  \draw[-latex, thick] (\xR-0.3,\yD-0.32) -- ++(0,-0.65) coordinate (ID2);
  \draw[-latex, thick] (0.55,-0.32) -- ++(0,-0.8) coordinate (IS);

  % Rótulos como nós explícitos, com folga de pelo menos 5 pt até fios, componentes e setas.
  % As setas de V+ e V- ocupam 0,42 cm acima/abaixo do nó vcc/vee; o zigue-zague de R
  % tem 0,21 cm de meia largura.
  \node[anchor=south, inner sep=0pt] at (0,\yT+0.4+0.42+0.27) {$V^+ = 5\,\mathrm{V}$};
  \node[anchor=north, inner sep=0pt] at (0,\yB-0.42-0.2) {$V^- = -5\,\mathrm{V}$};
  \node[anchor=east, align=left, inner sep=0pt] at (\xL-0.21-0.22,{(\yD+\yT)/2}) {$R_D$\\$25\,\mathrm{k}\Omega$};
  \node[anchor=west, align=left, inner sep=0pt] at (\xR+0.21+0.22,{(\yD+\yT)/2}) {$R_D$\\$25\,\mathrm{k}\Omega$};
  \node[anchor=east, align=left, inner sep=0pt] at (-0.21-0.22,\yB/2) {$R_S$\\$20\,\mathrm{k}\Omega$};
  \node[anchor=east, inner sep=0pt] at (\xL-0.23,\yD) {$D_1$};
  \node[anchor=north west, inner sep=0pt] at (\xR+0.22,\yD-0.22) {$D_2$};
  \node[anchor=west, inner sep=0pt] at (\xO+0.24,\yD) {$v_{O2}$};
  \node[anchor=east, inner sep=0pt] at ($(V1)+(-0.24,0)$) {$v_1$};
  \node[anchor=west, inner sep=0pt] at ($(V2)+(0.24,0)$) {$v_2$};
  \node[anchor=west, inner sep=0pt] at ($(ID1)+(0.25,0.325)$) {$I_{D1}$};
  \node[anchor=east, inner sep=0pt] at ($(ID2)+(-0.25,0.325)$) {$I_{D2}$};
  \node[anchor=west, inner sep=0pt] at ($(IS)+(0.25,0.4)$) {$I_S$};
  \node[anchor=west, inner sep=0pt] at ($(M1.gate -| M1.drain)+(0.33,0)$) {$M_1$};
  \node[anchor=east, inner sep=0pt] at ($(M2.gate -| M2.drain)+(-0.33,0)$) {$M_2$};
  \node[anchor=south, inner sep=0pt] at (0,0.32) {$S$};
\end{circuitikz}
\end{document}
